% Solution to Problem 3 — A Markov chain for interpolation Macdonald polynomials (Williams)
% v8, corrected against First Proof paper
% (Ben Dali-Williams, "A probabilistic interpretation for interpolation Macdonald polynomials")

\documentclass[11pt]{article}

\usepackage[margin=1in]{geometry}
\usepackage{amsmath, amssymb, amsthm}
\usepackage{mathtools}
\usepackage{xcolor}
\usepackage{hyperref}
\hypersetup{colorlinks=true,
            linkcolor=blue!50!black,
            citecolor=blue!50!black,
            urlcolor=blue!50!black}

\newcommand{\R}{\mathbb{R}}
\newcommand{\N}{\mathbb{N}}
\newcommand{\Z}{\mathbb{Z}}
\newcommand{\Snl}{S_n(\lambda)}
\newcommand{\wt}{\operatorname{wt}}

\newtheorem{theorem}{Theorem}
\newtheorem{lemma}{Lemma}
\newtheorem{proposition}{Proposition}
\newtheorem{corollary}{Corollary}
\theoremstyle{definition}
\newtheorem{definition}{Definition}
\newtheorem{remark}{Remark}

\title{Solution to Problem 3: A Probabilistic Interpretation\\
for Interpolation Macdonald Polynomials}
\author{}
\date{May 27, 2026}

\begin{document}
\maketitle

\section*{Statement and Answer}

Let $\lambda=(\lambda_1>\cdots>\lambda_n\ge0)$ be a partition with distinct
parts that is moreover \emph{restricted}: it has a unique part of size $0$ and no
part of size $1$. Let $\Snl$ be the orbit of $\lambda$ under permutation of
parts (equivalently, the compositions that rearrange $\lambda$). Let
$F^*_\mu(x_1,\dots,x_n;q,t)$ and $P^*_\lambda(x_1,\dots,x_n;q,t)$ be the
interpolation ASEP polynomial and interpolation Macdonald polynomial.

\medskip
\noindent\textbf{Theorem.} \emph{There is a nontrivial Markov chain on $\Snl$ ---
the \emph{interpolation $t$-Push TASEP} --- whose stationary distribution is}
\[
\pi^*_\lambda(\mu)\;=\;\frac{F^*_\mu(x_1,\dots,x_n;\,q=1,\,t)}{P^*_\lambda(x_1,\dots,x_n;\,q=1,\,t)},
\qquad \mu\in\Snl,
\]
\emph{and whose transition probabilities are not described using the
polynomials $F^*_\mu$.}

\medskip
\noindent The answer is \textbf{YES}. We caution at the outset that the chain is
\emph{not} the ordinary multispecies ASEP: that chain has the ordinary (not
interpolation) Macdonald polynomials as its stationary weights (Ayyer--Martin--Williams),
and substituting it here would solve a different, already-known problem. The
correct object is a $q=1$ \emph{push}-TASEP with an interpolation correction
encoded by a ringing ``bell''.

\section*{Probabilities and notation}

Throughout we work at $q=1$ and write $P^*_\lambda=P^*_\lambda(x;1,t)$, etc.
For $1\le k\le n$ put
\[
p_k:=\frac{t^{-n+1}(1-t)}{x_k-t^{-n+2}},\qquad
q_k:=\frac{(1-t)\,x_k}{x_k-t^{-n+2}} .
\]
For $0<t<1$ and $x_i>t^{-n+1}$ these lie in $[0,1]$ and serve as
skipping/displacement probabilities. We use $e^*_k(x;t)$ for the interpolation
elementary symmetric polynomials and the factorization (valid at $q=1$
\cite{Dolega,BDW25a})
\begin{equation}\label{eq:factor}
P^*_\lambda(x;1,t)=\prod_{1\le i\le\lambda_1}e^*_{\lambda'_i}(x;t),
\end{equation}
$\lambda'$ the conjugate partition.

\section*{The interpolation $t$-Push TASEP}

\begin{definition}\label{def:chain}
Fix $\lambda$ with $\lambda_n=0$. States of the chain are configurations of
particles on a ring of $n$ sites labelled by the parts of $\lambda$: state $\eta$
places the particle labelled $\eta_j$ at site $j$ (label $0$ is a vacancy). One
step from $\eta$ proceeds in three stages.

\smallskip
\noindent\textbf{(Step 0) The bell.} The bell at site $j$ rings with probability
\[
P_j=\frac{\prod_{k<j}\bigl(x_k-t^{-n+2}\bigr)\,\prod_{k>j}\bigl(x_k-t^{-n+1}\bigr)}
{e^*_{n-1}(x;t)},
\quad
e^*_{n-1}(x;t)=\sum_{j=1}^n\prod_{k<j}\bigl(x_k-t^{-n+2}\bigr)\prod_{k>j}\bigl(x_k-t^{-n+1}\bigr).
\]

\smallskip
\noindent\textbf{(Step 1) Push-TASEP move.} The particle at site $j$, say with
label $a$, is activated and travels clockwise by the $t$-Push TASEP rule: if
there are $m$ weaker particles in the system (labels $<a$, vacancies counting as
label $0$), the active particle moves to the $k$-th such weaker particle with
probability $t^{k-1}/[m]_t$. If that site held a positive label, that particle
becomes active and repeats; the cascade ends when the active particle reaches a
vacancy. Afterwards site $j$ is vacant; we treat this vacancy as a new particle
labelled $a:=0$.

\smallskip
\noindent\textbf{(Step 2) Vacancy sweep.} The label-$0$ particle goes to site $1$
and travels clockwise. Reaching a site $k$ ($1\le k\le j-1$) holding a label
$b\ge0$, it skips that site with probability $1-p_k$ if $b\ge a$ and $1-q_k$ if
$b<a$; otherwise it settles there, activating/displacing the resident, which then
continues clockwise toward site $j$ by the same rule. The active particle stops
once it displaces another or arrives at site $j$, where it settles.
\end{definition}

We denote the transition probabilities of Steps 1 and 2 by $P^{(1)}_j$ and
$P^{(2)}_j$, so that for $\mu,\nu\in\Snl$,
\[
P(\mu,\nu)=\sum_{1\le j\le n}P_j\sum_{\rho\in\Snl:\rho_j=0}
P^{(1)}_j(\mu,\rho)\,P^{(2)}_j(\rho,\nu).
\]

\begin{remark}[nontriviality]
The rates $P_j,p_k,q_k$ are explicit rational functions of $x_1,\dots,x_n$ and
$t$ only. They do \emph{not} reference $F^*_\mu$, so the chain is nontrivial in
the required sense (unlike a Metropolis--Hastings chain built from the target
distribution).
\end{remark}

\begin{theorem}\label{thm:main}
In the interpolation $t$-Push TASEP with content $\lambda$, the stationary
probability of $\mu\in\Snl$ is $\pi^*_\lambda(\mu)=F^*_\mu(x;1,t)/P^*_\lambda(x;1,t)$.
\end{theorem}

\section*{Two-line queues and the proof}

We use classical two-line queues \cite{CMW22} and signed two-line queues
\cite{BDW25a}, specialized to $q=1$. Let $\mathcal Q^\eta_\kappa$ be the classical
two-line queues with top row $\eta$ and bottom row $\kappa$, with pair-weight
generating function $a^\eta_\kappa=\sum_{Q\in\mathcal Q^\eta_\kappa}\wt_{\mathrm{pair}}(Q)$.
Let $\mathcal G^\alpha_\mu$ be the signed two-line queues with top row $\alpha\in\Z^n$
and bottom row $\mu$, with $b^\alpha_\mu=\sum_{Q\in\mathcal G^\alpha_\mu}\wt_{\mathrm{pair}}(Q)$,
and full weight $\wt(Q)=\wt_{\mathrm{pair}}(Q)\wt_{\mathrm{ball}}(Q)$ where
\begin{equation}\label{eq:wta}
\wt_\alpha\,b^\alpha_\mu=\sum_{Q\in\mathcal G^\alpha_\mu}\wt(Q),
\quad \wt_\alpha:=\prod_{k:\alpha_k>0}x_k\prod_{k:\alpha_k<0}\bigl(-t^{-(n-1)}\bigr).
\end{equation}

\begin{definition}\label{def:unsigned}
For a signed two-line queue $Q\in\mathcal G^\alpha_\mu$ let $\overline Q$ be its
unsigned version (forget signs on the top row); its top composition is
$\|\alpha\|=(|\alpha_1|,\dots,|\alpha_n|)$. Let $\mathcal G^\kappa_\mu$ be the set of
unsigned objects so obtained, with $\|\alpha\|=\kappa$. For $\overline Q\in\mathcal G^\kappa_\mu$
assign a nontrivial pairing $p$ the weight $(1-t)t^{\mathrm{skip}(p)}$, and a top-row
ball $B$ in column $k$ labelled $a>0$ with the ball below labelled $b$ the weight
\[
\wt(B)=\begin{cases}x_k-t^{-(n-1)} & b=a,\\ x_k & b>a,\\ t^{-(n-1)} & b<a.\end{cases}
\]
\end{definition}

\begin{lemma}\label{lem:signforget}
For $\lambda$ with distinct parts and $\kappa,\mu\in\Snl$, and $\overline Q\in\mathcal G^\kappa_\mu$,
$\wt(\overline Q)=\sum_Q\wt(Q)$, the sum over signed $Q$ forgetting to $\overline Q$.
\end{lemma}

\begin{proof}
Enumerate the admissible sign choices for each top-row ball $B$ labelled $a>0$:
a $-$ sign is forced when the ball below is a vacancy or labelled $b<a$; a $+$
sign is forced when $b>a$; both are allowed when $b=a$. One checks this matches
\eqref{eq:wta}: a $-$ sign multiplies the ball weight by $-1$ in passing from
$Q$ to $\overline Q$, but the weight of the pairing attached to $B$ is also
multiplied by $-1$, so signs cancel consistently.
\end{proof}

For $\kappa\in\Snl$ set $c^\kappa_\mu:=\sum_{\alpha:\|\alpha\|=\kappa}\wt_\alpha\,b^\alpha_\mu$.
Combining \eqref{eq:wta} and Lemma~\ref{lem:signforget}:

\begin{lemma}\label{lem:ckm}
$c^\kappa_\mu=\sum_{\overline Q\in\mathcal G^\kappa_\mu}\wt(\overline Q)$. Since $\lambda$ has distinct
parts, $\mathcal G^\kappa_\mu$ is empty or a single element.
\end{lemma}

\subsection*{The closing recursion (where the hypotheses on $\lambda$ enter)}

Fix a weakly order-preserving $\phi:\N\to\N$ and partitions $\lambda,\kappa$ with
$\phi(\lambda)=\kappa$. For $\eta\in S_n(\kappa)$ set
$G^*_\eta(x;t):=\sum_{\rho\in\Snl:\,\phi(\rho)=\eta}F^*_\rho(x;1,t)$. The following
is an interpolation analogue of a result of Ayyer--Martin--Williams
\cite{AMW25}, provable in the same way using \cite{AS19}.

\begin{theorem}\label{thm:G}
For $\lambda,\kappa$ as above and all $\eta\in S_n(\kappa)$, at $q=1$,
\[
\frac{G^*_\eta(x;t)}{P^*_\lambda(x;1,t)}=\frac{F^*_\eta(x;1,t)}{P^*_\kappa(x;1,t)} .
\]
\end{theorem}

\begin{corollary}\label{cor:rec}
Let $\rho$ be a composition with $\rho_i\ne1$ for all $i$, with $k$ nonzero
parts, and set $\eta=\rho^-$ (where $\rho^-_i=\max(\rho_i-1,0)$). Then at $q=1$,
\[
F^*_\rho(x;1,t)=F^*_\eta(x;1,t)\cdot e^*_k(x;t).
\]
\end{corollary}

\begin{proof}
Let $\lambda,\kappa$ reorder $\rho,\eta$. Take $\phi:i\mapsto\max(i-1,0)$, so
$\phi(\rho)=\eta$. \emph{Because $\lambda$ has no part of size $1$}, $\phi$ is a
bijection from $\{0,2,3,\dots\}$ onto $\{0,1,2,\dots\}$, so $\rho$ is the unique
composition in $\Snl$ with $\phi(\rho)=\eta$; hence $G^*_\eta=F^*_\rho$. By
Theorem~\ref{thm:G},
$F^*_\rho/P^*_\lambda=F^*_\eta/P^*_\kappa$. Finally, by \eqref{eq:factor} and the
fact that $\kappa$ is obtained from $\lambda$ by deleting the largest column (of
size $k$), $P^*_\lambda/P^*_\kappa=e^*_k(x;t)$, giving the claim.
\end{proof}

This is exactly where the restriction on $\lambda$ is essential: \emph{no part of
size $1$} makes $\phi$ bijective so the recursion closes, and the \emph{unique
part of size $0$} guarantees $\eta$ has a unique $0$-part so the column-deletion
is unambiguous. The hypotheses are structural, not (as one might guess) a device
to avoid a $0/0$ normalization.

\subsection*{From two-line queues to the dynamics}

\begin{proposition}\label{prop:step2}
For $\rho,\nu\in\Snl$ with $\rho_j=0$,
\[
P^{(2)}_j(\rho,\nu)=\frac{c^\rho_\nu}{\prod_{k<j}(x_k-t^{-n+2})\prod_{k>j}(x_k-t^{-n+1})},
\quad\text{equivalently}\quad
P_j\,P^{(2)}_j(\rho,\nu)=\frac{c^\rho_\nu}{e^*_{n-1}} .
\]
\end{proposition}

\begin{proof}[Idea]
Step 2 is encoded by the unique element of $\mathcal G^\rho_\nu$ (Lemma~\ref{lem:ckm}):
a particle labelled $a$ moving from column $k$ to $k'$ corresponds to a
nontrivial pairing joining a top ball in column $k$ to a bottom ball in column
$k'$; stationary particles give trivial pairings. Dividing each ball/pairing
weight in $\wt(\overline Q)$ by the matching factor of
$D=\prod_{k<j}(x_k-t^{-n+2})\prod_{k>j}(x_k-t^{-n+1})$ reproduces, case by case,
the skip/settle probabilities $1-p_k,\,1-q_k,\,p_k,\,q_k$ of Step 2 (and a
factor $1$ for $k>j$, which never moves). Hence $\wt(\overline Q)/D=P^{(2)}_j(\rho,\nu)$.
\end{proof}

\begin{proposition}\label{prop:transition}
If $\lambda$ is restricted and $\mu,\nu\in\Snl$, then
$\displaystyle P(\mu,\nu)=\sum_{\rho\in\Snl}\frac{a^\mu_\rho\,c^\rho_\nu}{e^*_{n-1}}$.
\end{proposition}

\begin{proof}
Combine the encoding of Step 1 by classical two-line queues
\cite[Lem.~5.4]{AMW25} (giving $P^{(1)}_j(\mu,\rho)$ via $a^\mu_\rho$) with
Proposition~\ref{prop:step2}:
\[
P(\mu,\nu)=\sum_j P_j\!\!\sum_{\rho:\rho_j=0}\!P^{(1)}_j(\mu,\rho)P^{(2)}_j(\rho,\nu)
=\sum_j\sum_{\rho:\rho_j=0}\frac{a^\mu_\rho c^\rho_\nu}{e^*_{n-1}}
=\sum_{\rho\in\Snl}\frac{a^\mu_\rho c^\rho_\nu}{e^*_{n-1}}.
\]
\end{proof}

\subsection*{Proof of Theorem~\ref{thm:main}}

\begin{proof}
Fix $\nu\in\Snl$. By \cite[Thm.~1.15, Lem.~5.6]{BDW25a},
\[
F^*_\nu(x;1,t)=\sum_{\eta\in\N^n}F^{*\eta}_\nu(x;t)\,F^*_{\eta^-}(x;1,t),
\qquad
F^{*\eta}_\nu(x;t)=\sum_{\kappa\in\N^n}a^\eta_\kappa\,c^\kappa_\nu .
\]
By Corollary~\ref{cor:rec} (using that $\eta$ has a unique $0$-part),
$F^*_{\eta^-}(x;1,t)=F^*_\eta(x;1,t)/e^*_{n-1}(x;t)$. Substituting,
\[
F^*_\nu(x;1,t)=\sum_{\eta\in\N^n}F^*_\eta(x;1,t)\sum_{\kappa\in\N^n}
\frac{a^\eta_\kappa\,c^\kappa_\nu}{e^*_{n-1}(x;t)}
=\sum_{\eta\in\N^n}F^*_\eta(x;1,t)\,P(\eta,\nu),
\]
the last equality by Proposition~\ref{prop:transition}. Thus the vector
$\bigl(F^*_\mu(x;1,t)\bigr)_\mu$ is left-fixed by the transition matrix, i.e.\
proportional to the stationary distribution. Since
$P^*_\lambda=\sum_{\mu\in\Snl}F^*_\mu$, normalizing gives
$\pi^*_\lambda(\mu)=F^*_\mu(x;1,t)/P^*_\lambda(x;1,t)$.
\end{proof}

\section*{Remark on the failure mode to avoid}

Replacing this chain by the ordinary multispecies ASEP (matrix-product ansatz /
Zamolodchikov--Faddeev algebra) does \emph{not} solve the stated problem: that
construction yields the ordinary Macdonald polynomials as stationary weights
\cite{CMW22,AMW25}, not the \emph{interpolation} polynomials $F^*_\mu,P^*_\lambda$.
The interpolation correction is precisely what the bell $P_j$ and the
push-with-sweep dynamics supply, and the proof above runs through signed
two-line queues rather than a trace over an auxiliary algebra.

\begin{thebibliography}{9}
\bibitem{AS19} P.~Alexandersson, M.~Sawhney.
  \emph{Properties of non-symmetric Macdonald polynomials at $q=1$ and $q=0$}.
  Ann.\ Comb.\ \textbf{23} (2019), 219--239.
\bibitem{AMW25} A.~Ayyer, J.~Martin, L.~Williams.
  \emph{The inhomogeneous $t$-PushTASEP and Macdonald polynomials at $q=1$}.
  Ann.\ Inst.\ Henri Poincar\'e D, 2025.
\bibitem{BDW25a} H.~Ben Dali, L.~Williams.
  \emph{A combinatorial formula for interpolation Macdonald polynomials}.
  Preprint, arXiv:2510.02587, 2025.
\bibitem{BDW26} H.~Ben Dali, L.~Williams.
  \emph{A probabilistic interpretation for interpolation Macdonald polynomials}.
  Preprint, arXiv:2602.13492, 2026.
\bibitem{CMW22} S.~Corteel, O.~Mandelshtam, L.~Williams.
  \emph{From multiline queues to Macdonald polynomials via the exclusion process}.
  Amer.\ J.\ Math.\ \textbf{144} (2022), 395--436.
\bibitem{Dolega} M.~Dolega.
  \emph{Strong factorization property of Macdonald polynomials and higher-order
  Macdonald's positivity conjecture}.
  J.\ Algebraic Combin.\ \textbf{46} (2017), 135--163.
\end{thebibliography}

\end{document}